\documentclass[11pt,a4paper]{article}
\usepackage[margin=25mm]{geometry}
\usepackage{amsmath,amssymb,amsthm,hyperref}
\setlength{\emergencystretch}{2em}
\newtheorem{lemma}{Lemma}
\newtheorem{theorem}[lemma]{Theorem}
\newtheorem{corollary}[lemma]{Corollary}
\newcommand{\gen}[1]{\langle#1\rangle}
\newcommand{\Sym}{\operatorname{Sym}}
\newcommand{\Alt}{\operatorname{Alt}}
\title{Odd permutations at an attachment and the remaining kernel obstruction}
\author{Human Compression project: a component of the L019 recovery}
\date{}
\begin{document}\maketitle
For a set $X$, write $\Sym(X)$ for permutations supported on $X$,
fixing every other point. We use actual supported subgroups throughout.
The acquired leaf criterion says that for an anchored string C-marking
$(h_0,\ldots,h_{k-1})$, where all $h_j$ with $j>0$ fix $a$,
adjoining a pure new transposition at $a$ preserves IP exactly when
\[
 P_j\cap\Sym(X_{j-1})\le P_{j-1}\quad(1\le j\le k),
 \qquad P_j=\gen{h_0,\ldots,h_{j-1}},\quad X_j=a^{P_j}.
                                                               \tag{1}
\]
The first condition is automatic. We prove a uniform way to obtain
these conditions from the IP of a larger, already given marking.

\begin{lemma}\label{overlap}
Let groups $P,E$ have supports contained respectively in sets
$\Omega,Y$ with $\Omega\cap Y=\{a\}$. Suppose $E$ is transitive
on $Y$, $|Y|\ge2$, and put $X=a^P$. If $|X|\ge2$, then
\[
 \Alt(X\cup Y)\le\gen{P,E}.
\]
If $E$ contains an odd permutation, then
$\Sym(X\cup Y)\le\gen{P,E}$.
\end{lemma}
\begin{proof}
For any $x\in X\setminus\{a\}$ and $y\in Y\setminus\{a\}$,
choose $p\in P,e\in E$ with $a^p=x,a^e=y$.
The supports of $p,e$ intersect only at $a$. A direct evaluation
in the right-action convention shows that
$p^{-1}e^{-1}pe=(a\ y\ x)$, fixing every other point.
Thus the generated group contains every three-cycle with $a$,
one other point from $X$, and one from $Y$.

These cycles generate $\Alt(X\cup Y)$. To see this directly,
fix $x_0\in X\setminus\{a\}$ and $y_0\in Y\setminus\{a\}$.
For two distinct points $x,x'\in X\setminus\{a\}$, the product
$(a\ y_0\ x)^{-1}(a\ y_0\ x')$ is the three-cycle
$(a\ x\ x')$.
The same argument with $x_0$ produces every three-cycle on
$a,y,y'$ for distinct $y,y'\in Y\setminus\{a\}$.
Together with the mixed cycles, these are all three-cycles containing
$a$. They generate the alternating group: any three-cycle avoiding
$a$ is a product of two such cycles, and every even permutation is
a product of three-cycles. The argument also covers $|X|=|Y|=2$.
All displayed cycles are supported on $X\cup Y$.

An odd element of $E$ is supported on $Y\subseteq X\cup Y$.
Adjoining it to the alternating group gives the full supported
symmetric group on that union, proving the final assertion.
\end{proof}

\begin{theorem}\label{kernels}
Suppose $(h_0,\ldots,h_{k-1})$ is a submarking of a string
C-marking and has endpoint anchor $a$. Let $E$ be generated by
a disjoint set of labels from the same C-marking. Suppose $E$ is
transitive on a set $Y$ of size at least two and is supported there,
while all the $h_j$ are supported in $\Omega$ with
$\Omega\cap Y=\{a\}$. Then every even part of a kernel in (1)
is already in the preceding prefix:
\[
 P_j\cap\Alt(X_{j-1})\le P_{j-1}.                    \tag{2}
\]
If $E$ contains an odd permutation, all conditions (1) hold and
the pure-leaf extension is C.

For arbitrary $E$, a failed condition (1) has precisely an index-two
sign defect. In detail, with
\[
 Q_j=P_j\cap\Sym(X_{j-1}),\qquad
 K_j=P_{j-1}\cap\Sym(X_{j-1}),
\]
one either has $Q_j=K_j$, or
$K_j=Q_j\cap\Alt(X_{j-1})$ and $[Q_j:K_j]=2$.
\end{theorem}
\begin{proof}
For $j=1$, the supported symmetric group on $X_0=\{a\}$
is trivial. For $j\ge2$, the anchor is moved by $h_0$, so
$|X_{j-1}|\ge2$. Apply Lemma~\ref{overlap} with $P=P_{j-1}$.
It puts the supported $\Alt(X_{j-1})$, and also
$\Sym(X_{j-1})$ when $E$ has an odd element, inside
$\gen{P_{j-1},E}$.
The original intersection property gives the exact bound
\[
 P_j\cap\gen{P_{j-1},E}=P_{j-1},
\]
since the label sets of $E$ and the $h_j$ are disjoint.
This proves (2) and, in the odd case, every condition (1).
The complete acquired leaf theorem then proves the C-property of
the enlarged marking, including all its subgroup intersections.

Finally $K_j\le Q_j$, and (2) gives
$Q_j\cap\Alt(X_{j-1})\le K_j$.
If $K_j$ has an odd element it equals $Q_j$, because its even
subgroup already contains the entire even part of $Q_j$.
Otherwise it equals that even part. A strict inclusion is therefore
possible only when $Q_j$ has odd elements and $K_j$ does not,
in which case its index is exactly two.
\end{proof}

\begin{corollary}\label{split}
Let an internal perfect split of a transitive symmetric string
C-marking have the form
\[
 q=\rho_i=(a\ b)\beta,
 \quad L=\gen{\rho_0,\ldots,\rho_{i-1}},\quad
 J=\gen{q,\rho_{i+1},\ldots,\rho_{r-1}}.
\]
Here the old left internal factor is trivial. If both $L$ and $J$
contain odd permutations, subdividing the crossing by replacing
$q$ with $(a\ c),(c\ b)\beta$ at a new point $c$ gives a string
C-marking of degree $n+1$ and rank $r+1$, generating $S_{n+1}$.
The reversed assertion holds when the right internal factor is trivial.

Without the two oddness assumptions, at most one of the two leaf
families in the acquired split criterion can fail, and every failure
there is exactly the sign defect of Theorem~\ref{kernels}.
\end{corollary}
\begin{proof}
The supported left group $L$ is transitive on $A$. Since the split
is internal, $|A|\ge2$. The group $J$ is supported and transitive
on $\{a\}\cup B$: its tail is transitive on $B$, and $q$
joins $a$ to $b$. These two domains meet only at $a$.
Apply Theorem~\ref{kernels} first to the old left wing in reversed
order, with external group $J$, and then to the marking beginning
with $q$ and continuing rightwards, with external group $L$.
These are exactly the two original lists in the repaired split
criterion. Their endpoint-anchor hypotheses follow from the unique
crossing and the actual wing transitivity. If both external groups
have odd elements, both lists satisfy all the required tests.
The complete split criterion gives the full C-property and string
relations of the subdivided marking.

For its group, write $u=(a\ c),v=(c\ b)\beta$.
Then $uvu=q$ on the old domain and fixes $c$.
All other old marks are retained, so the new group contains the old
$S_n$ fixing $c$, together with $u$. Conjugating $u$ by that $S_n$
gives all transpositions $(x\ c)$; hence the group is $S_{n+1}$.
The inserted involutions are distinct and nonidentity by the split
criterion, so the rank increases by one.

Lastly $\gen{L,J}=S_n$, so at least one of these two actual groups
contains an odd permutation. Its external use proves one whole leaf
family. The other family has only the exact possible index-two defects
already established; no assertion is made that these remaining defects
are automatically absent in the high-rank range.
\end{proof}

\paragraph{Attribution and boundary.}
The extension and its intended application are those of
Cameron--Fernandes--Leemans, \emph{The number of string C-groups
of high rank}, arXiv:2212.12723v2. The complete restriction-kernel
criterion is the separately reviewed repair retained as an acquired
premise, including its standard C-group end criterion attribution.
The elementary commutator argument above supplies a new sufficient
interface from the original IP, and identifies the only residual
obstruction when the attaching group is even. No priority claim is
made for the alternating-group generation identity itself.

This proves extension for the stated oddness condition. Existence
of such a split for every high-rank marking, or elimination of the
remaining sign defects by another method, has not been established.
The reversible correspondence, exact base constants and whole L019
remain pending. No new census, Lean or compression measurement enters.
\end{document}
